% Curs 10 - Pagina 9 (6-XII-1994).
% Fragment pentru % Curs 10 - Pagina 9 (6-XII-1994).
% Fragment pentru % Curs 10 - Pagina 9 (6-XII-1994).
% Fragment pentru % Curs 10 - Pagina 9 (6-XII-1994).
% Fragment pentru \input{courses/pages/curs10-pagina9-ro}.
% Fara preambul si fara intrerupere fortata de pagina.
% Continuarea demonstratiei T. lui Sylow 1 de pe pagina 8.
% Note editoriale: pasajele incerte sau logic problematice sunt semnalate
% prin comentarii sau note, nu corectate tacit.

\subsection*{Curs 10 - Pagina 9}

\noindent $\varphi(K)=\dfrac{K}{H}=K'.$

\medskip
\noindent Iar $H\trianglelefteq K$ pentru că $H\trianglelefteq N_G(H)$.

\medskip
\noindent Aplicăm Lagrange:
\[
|K|=|H|\cdot[K:H]=|H|\cdot[\,\ldots\,]
\quad\Longrightarrow\quad p\mid|K'|.
\]
% In manuscris, termenul dintre ultimele paranteze patrate este omis/neclar.
% La capatul randului pare scris „p | |K'|”; argumentul matematic
% ar trebui sa stabileasca |K|=p^{i+1}. Nu inlocuim lectura originala.

\bigskip
\noindent\fbox{\textbf{Dem. T. lui Sylow 2}}

\medskip
\noindent Fie $P\in\operatorname{Syl}_p(G)$. Fie
\[
A=\{aP\mid a\in G\}.
\]

\noindent $H\times A\longrightarrow A$,
\[
(h,aP)\longmapsto(ha)P\in A,
\]
\noindent acţiune prin translaţie.

\noindent Verifică că este o acţiune.

\medskip
\noindent Considerăm $A_0$ corespunzător:
\[
A_0=\{aP\mid haP=aP,\ (\forall)\,h\in H\}.
\]

\noindent $A_0\neq\varnothing$, fiindcă $p\nmid[G:P]$ şi
$|P|=p^t$.

\medskip
\noindent Din lema 2 $\Longrightarrow$
\[
p\nmid[N_G(P):P].
\]
% Lectura „[N_G(P):P]” este vizibila, dar aplicarea lemei la
% actiunea lui H pe G/P ar cere mai intai concluzia p\nmid|A_0|.
% Necesita verificare in manuscris; nu s-a corectat tacit.

% Linia „Aplic. lema 2” este taiata in manuscris.
\medskip
\noindent
\[
|A_0|\equiv|A|\pmod p
\quad\Longrightarrow\quad p\nmid|A_0|
\quad\Longrightarrow\quad A_0\neq\varnothing.
\]
\noindent $|A|=m$ şi $p\nmid m$.

\medskip
\noindent $(\exists)\,aP\in A_0$. Atunci
\[
haP=aP,\quad(\forall)\,h\in H
\quad\Longleftrightarrow\quad
 a^{-1}haP=P,\quad(\forall)\,h\in H.
\]
\[
\Longleftrightarrow\quad a^{-1}Ha\leq P
\quad\Longrightarrow\quad H\leq aPa^{-1}.
\]

\medskip
\noindent Fie $Q\in\operatorname{Syl}_p(G)$. Atunci
\[
aQa^{-1}\leq P.
\]
\noindent Dar
\[
|P|=p^t,\qquad |Q|=p^t
\quad\Longrightarrow\quad |aQa^{-1}|=p^t
\quad\Longrightarrow\quad aQa^{-1}=P.
\]

\bigskip
\noindent\underline{Aplicaţie}: $|G|=pq$, $p>q$, $q\nmid(p-1)$.

\medskip
\noindent Cauchy: $(\exists)\,a,b\in G$ cu
$o(a)=p$, $o(b)=q$.

\noindent $P=\langle a\rangle$, $Q=\langle b\rangle$.

\noindent Din T. 3:
\[
|\operatorname{Syl}_p(G)|\mid pq,
\qquad |\operatorname{Syl}_p(G)|\equiv1\pmod p
\quad\Longrightarrow\quad |\operatorname{Syl}_p(G)|=1.
\]

\noindent $P\in\operatorname{Syl}_p(G)\ \ni Q$.%
% Lectura finala „\ni Q” este nesigura si, interpretata ca
% incluziune Q<=P, ar fi incompatibila cu |Q|=q si |P|=p.
% Pasajul se incheie aici pe pagina 9.
.
% Fara preambul si fara intrerupere fortata de pagina.
% Continuarea demonstratiei T. lui Sylow 1 de pe pagina 8.
% Note editoriale: pasajele incerte sau logic problematice sunt semnalate
% prin comentarii sau note, nu corectate tacit.

\subsection*{Curs 10 - Pagina 9}

\noindent $\varphi(K)=\dfrac{K}{H}=K'.$

\medskip
\noindent Iar $H\trianglelefteq K$ pentru că $H\trianglelefteq N_G(H)$.

\medskip
\noindent Aplicăm Lagrange:
\[
|K|=|H|\cdot[K:H]=|H|\cdot[\,\ldots\,]
\quad\Longrightarrow\quad p\mid|K'|.
\]
% In manuscris, termenul dintre ultimele paranteze patrate este omis/neclar.
% La capatul randului pare scris „p | |K'|”; argumentul matematic
% ar trebui sa stabileasca |K|=p^{i+1}. Nu inlocuim lectura originala.

\bigskip
\noindent\fbox{\textbf{Dem. T. lui Sylow 2}}

\medskip
\noindent Fie $P\in\operatorname{Syl}_p(G)$. Fie
\[
A=\{aP\mid a\in G\}.
\]

\noindent $H\times A\longrightarrow A$,
\[
(h,aP)\longmapsto(ha)P\in A,
\]
\noindent acţiune prin translaţie.

\noindent Verifică că este o acţiune.

\medskip
\noindent Considerăm $A_0$ corespunzător:
\[
A_0=\{aP\mid haP=aP,\ (\forall)\,h\in H\}.
\]

\noindent $A_0\neq\varnothing$, fiindcă $p\nmid[G:P]$ şi
$|P|=p^t$.

\medskip
\noindent Din lema 2 $\Longrightarrow$
\[
p\nmid[N_G(P):P].
\]
% Lectura „[N_G(P):P]” este vizibila, dar aplicarea lemei la
% actiunea lui H pe G/P ar cere mai intai concluzia p\nmid|A_0|.
% Necesita verificare in manuscris; nu s-a corectat tacit.

% Linia „Aplic. lema 2” este taiata in manuscris.
\medskip
\noindent
\[
|A_0|\equiv|A|\pmod p
\quad\Longrightarrow\quad p\nmid|A_0|
\quad\Longrightarrow\quad A_0\neq\varnothing.
\]
\noindent $|A|=m$ şi $p\nmid m$.

\medskip
\noindent $(\exists)\,aP\in A_0$. Atunci
\[
haP=aP,\quad(\forall)\,h\in H
\quad\Longleftrightarrow\quad
 a^{-1}haP=P,\quad(\forall)\,h\in H.
\]
\[
\Longleftrightarrow\quad a^{-1}Ha\leq P
\quad\Longrightarrow\quad H\leq aPa^{-1}.
\]

\medskip
\noindent Fie $Q\in\operatorname{Syl}_p(G)$. Atunci
\[
aQa^{-1}\leq P.
\]
\noindent Dar
\[
|P|=p^t,\qquad |Q|=p^t
\quad\Longrightarrow\quad |aQa^{-1}|=p^t
\quad\Longrightarrow\quad aQa^{-1}=P.
\]

\bigskip
\noindent\underline{Aplicaţie}: $|G|=pq$, $p>q$, $q\nmid(p-1)$.

\medskip
\noindent Cauchy: $(\exists)\,a,b\in G$ cu
$o(a)=p$, $o(b)=q$.

\noindent $P=\langle a\rangle$, $Q=\langle b\rangle$.

\noindent Din T. 3:
\[
|\operatorname{Syl}_p(G)|\mid pq,
\qquad |\operatorname{Syl}_p(G)|\equiv1\pmod p
\quad\Longrightarrow\quad |\operatorname{Syl}_p(G)|=1.
\]

\noindent $P\in\operatorname{Syl}_p(G)\ \ni Q$.%
% Lectura finala „\ni Q” este nesigura si, interpretata ca
% incluziune Q<=P, ar fi incompatibila cu |Q|=q si |P|=p.
% Pasajul se incheie aici pe pagina 9.
.
% Fara preambul si fara intrerupere fortata de pagina.
% Continuarea demonstratiei T. lui Sylow 1 de pe pagina 8.
% Note editoriale: pasajele incerte sau logic problematice sunt semnalate
% prin comentarii sau note, nu corectate tacit.

\subsection*{Curs 10 - Pagina 9}

\noindent $\varphi(K)=\dfrac{K}{H}=K'.$

\medskip
\noindent Iar $H\trianglelefteq K$ pentru că $H\trianglelefteq N_G(H)$.

\medskip
\noindent Aplicăm Lagrange:
\[
|K|=|H|\cdot[K:H]=|H|\cdot[\,\ldots\,]
\quad\Longrightarrow\quad p\mid|K'|.
\]
% In manuscris, termenul dintre ultimele paranteze patrate este omis/neclar.
% La capatul randului pare scris „p | |K'|”; argumentul matematic
% ar trebui sa stabileasca |K|=p^{i+1}. Nu inlocuim lectura originala.

\bigskip
\noindent\fbox{\textbf{Dem. T. lui Sylow 2}}

\medskip
\noindent Fie $P\in\operatorname{Syl}_p(G)$. Fie
\[
A=\{aP\mid a\in G\}.
\]

\noindent $H\times A\longrightarrow A$,
\[
(h,aP)\longmapsto(ha)P\in A,
\]
\noindent acţiune prin translaţie.

\noindent Verifică că este o acţiune.

\medskip
\noindent Considerăm $A_0$ corespunzător:
\[
A_0=\{aP\mid haP=aP,\ (\forall)\,h\in H\}.
\]

\noindent $A_0\neq\varnothing$, fiindcă $p\nmid[G:P]$ şi
$|P|=p^t$.

\medskip
\noindent Din lema 2 $\Longrightarrow$
\[
p\nmid[N_G(P):P].
\]
% Lectura „[N_G(P):P]” este vizibila, dar aplicarea lemei la
% actiunea lui H pe G/P ar cere mai intai concluzia p\nmid|A_0|.
% Necesita verificare in manuscris; nu s-a corectat tacit.

% Linia „Aplic. lema 2” este taiata in manuscris.
\medskip
\noindent
\[
|A_0|\equiv|A|\pmod p
\quad\Longrightarrow\quad p\nmid|A_0|
\quad\Longrightarrow\quad A_0\neq\varnothing.
\]
\noindent $|A|=m$ şi $p\nmid m$.

\medskip
\noindent $(\exists)\,aP\in A_0$. Atunci
\[
haP=aP,\quad(\forall)\,h\in H
\quad\Longleftrightarrow\quad
 a^{-1}haP=P,\quad(\forall)\,h\in H.
\]
\[
\Longleftrightarrow\quad a^{-1}Ha\leq P
\quad\Longrightarrow\quad H\leq aPa^{-1}.
\]

\medskip
\noindent Fie $Q\in\operatorname{Syl}_p(G)$. Atunci
\[
aQa^{-1}\leq P.
\]
\noindent Dar
\[
|P|=p^t,\qquad |Q|=p^t
\quad\Longrightarrow\quad |aQa^{-1}|=p^t
\quad\Longrightarrow\quad aQa^{-1}=P.
\]

\bigskip
\noindent\underline{Aplicaţie}: $|G|=pq$, $p>q$, $q\nmid(p-1)$.

\medskip
\noindent Cauchy: $(\exists)\,a,b\in G$ cu
$o(a)=p$, $o(b)=q$.

\noindent $P=\langle a\rangle$, $Q=\langle b\rangle$.

\noindent Din T. 3:
\[
|\operatorname{Syl}_p(G)|\mid pq,
\qquad |\operatorname{Syl}_p(G)|\equiv1\pmod p
\quad\Longrightarrow\quad |\operatorname{Syl}_p(G)|=1.
\]

\noindent $P\in\operatorname{Syl}_p(G)\ \ni Q$.%
% Lectura finala „\ni Q” este nesigura si, interpretata ca
% incluziune Q<=P, ar fi incompatibila cu |Q|=q si |P|=p.
% Pasajul se incheie aici pe pagina 9.
.
% Fara preambul si fara intrerupere fortata de pagina.
% Continuarea demonstratiei T. lui Sylow 1 de pe pagina 8.
% Note editoriale: pasajele incerte sau logic problematice sunt semnalate
% prin comentarii sau note, nu corectate tacit.

\subsection*{Curs 10 - Pagina 9}

\noindent $\varphi(K)=\dfrac{K}{H}=K'.$

\medskip
\noindent Iar $H\trianglelefteq K$ pentru că $H\trianglelefteq N_G(H)$.

\medskip
\noindent Aplicăm Lagrange:
\[
|K|=|H|\cdot[K:H]=|H|\cdot[\,\ldots\,]
\quad\Longrightarrow\quad p\mid|K'|.
\]
% In manuscris, termenul dintre ultimele paranteze patrate este omis/neclar.
% La capatul randului pare scris „p | |K'|”; argumentul matematic
% ar trebui sa stabileasca |K|=p^{i+1}. Nu inlocuim lectura originala.

\bigskip
\noindent\fbox{\textbf{Dem. T. lui Sylow 2}}

\medskip
\noindent Fie $P\in\operatorname{Syl}_p(G)$. Fie
\[
A=\{aP\mid a\in G\}.
\]

\noindent $H\times A\longrightarrow A$,
\[
(h,aP)\longmapsto(ha)P\in A,
\]
\noindent acţiune prin translaţie.

\noindent Verifică că este o acţiune.

\medskip
\noindent Considerăm $A_0$ corespunzător:
\[
A_0=\{aP\mid haP=aP,\ (\forall)\,h\in H\}.
\]

\noindent $A_0\neq\varnothing$, fiindcă $p\nmid[G:P]$ şi
$|P|=p^t$.

\medskip
\noindent Din lema 2 $\Longrightarrow$
\[
p\nmid[N_G(P):P].
\]
% Lectura „[N_G(P):P]” este vizibila, dar aplicarea lemei la
% actiunea lui H pe G/P ar cere mai intai concluzia p\nmid|A_0|.
% Necesita verificare in manuscris; nu s-a corectat tacit.

% Linia „Aplic. lema 2” este taiata in manuscris.
\medskip
\noindent
\[
|A_0|\equiv|A|\pmod p
\quad\Longrightarrow\quad p\nmid|A_0|
\quad\Longrightarrow\quad A_0\neq\varnothing.
\]
\noindent $|A|=m$ şi $p\nmid m$.

\medskip
\noindent $(\exists)\,aP\in A_0$. Atunci
\[
haP=aP,\quad(\forall)\,h\in H
\quad\Longleftrightarrow\quad
 a^{-1}haP=P,\quad(\forall)\,h\in H.
\]
\[
\Longleftrightarrow\quad a^{-1}Ha\leq P
\quad\Longrightarrow\quad H\leq aPa^{-1}.
\]

\medskip
\noindent Fie $Q\in\operatorname{Syl}_p(G)$. Atunci
\[
aQa^{-1}\leq P.
\]
\noindent Dar
\[
|P|=p^t,\qquad |Q|=p^t
\quad\Longrightarrow\quad |aQa^{-1}|=p^t
\quad\Longrightarrow\quad aQa^{-1}=P.
\]

\bigskip
\noindent\underline{Aplicaţie}: $|G|=pq$, $p>q$, $q\nmid(p-1)$.

\medskip
\noindent Cauchy: $(\exists)\,a,b\in G$ cu
$o(a)=p$, $o(b)=q$.

\noindent $P=\langle a\rangle$, $Q=\langle b\rangle$.

\noindent Din T. 3:
\[
|\operatorname{Syl}_p(G)|\mid pq,
\qquad |\operatorname{Syl}_p(G)|\equiv1\pmod p
\quad\Longrightarrow\quad |\operatorname{Syl}_p(G)|=1.
\]

\noindent $P\in\operatorname{Syl}_p(G)\ \ni Q$.%
% Lectura finala „\ni Q” este nesigura si, interpretata ca
% incluziune Q<=P, ar fi incompatibila cu |Q|=q si |P|=p.
% Pasajul se incheie aici pe pagina 9.
